\documentclass[direction=tate]{article}
\usepackage[fontset=ipaex-type1]{pdftexja}
\usepackage{ruby}
\begin{document}
\section{縦組の見出し}\label{節:縦}
縦組の本文である。小仮名（ぁぃぅ）と長音ー、括弧「あ」を含む。
\tatechuyoko{12}と約物。参照は\ref{節:縦}。
$\mathmc{和文}+x_{\mathgt{添字}}=1$。
\ruby{漢字}{かんじ}。
\begin{pdftexjaparagraph}
縦組の厳密段落である。約物「括弧」と句読点、Latin を含む。
\end{pdftexjaparagraph}
本文の脚注\footnote{脚注の本文である。}。
\newpage
二ページ目の本文である。
\end{document}
